% WP4 supplementary atomic traceability matrix -- not included in main.tex by default
\section*{Phụ lục: Ma trận truy vết mệnh đề ràng buộc nguồn nguyên tử}
\begingroup\tiny\sloppy\setlength{\emergencystretch}{3em}\renewcommand{\arraystretch}{1.02}
\begin{longtable}{@{}p{0.9cm}p{3.0cm}p{0.55cm}p{4.05cm}p{2.0cm}p{3.45cm}@{}}
\toprule Mã & Nguồn / vị trí & Lớp & Mệnh đề nguồn & AE & Trạng thái / tiêu chí \\ \midrule\endfirsthead
\multicolumn{6}{c}{Ma trận truy vết -- tiếp theo}\\ \toprule Mã & Nguồn / vị trí & Lớp & Mệnh đề nguồn & AE & Trạng thái / tiêu chí \\ \midrule\endhead\bottomrule\endfoot
\texttt{C01.01} & QĐ 3090/QĐ-BKHCN; Điều 1 & A & Khung kiến trúc tổng thể quốc gia số là hệ quy chiếu tổng thể cho kiến trúc số. & AE01 & ĐÁP ỨNG: Hình định vị chỉ ra khung quốc gia và không tuyên bố kiến trúc miền thay thế khung quốc gia. \\ \addlinespace[1pt]
\texttt{C01.02} & QĐ 292/QĐ-BKHCN; Phụ lục, Chương 2, Mục II & A & Mô hình tham chiếu nghiệp vụ phải được xem xét khi mô tả miền nghiệp vụ. & AE01 & ĐÁP ỨNG: Có viewpoint/model tương ứng và truy vết tới quyết định kiến trúc. \\ \addlinespace[1pt]
\texttt{C01.03} & QĐ 292/QĐ-BKHCN; Phụ lục, Chương 2, Mục II & A & Mô hình tham chiếu dữ liệu phải được xem xét khi mô tả dữ liệu. & AE01 & ĐÁP ỨNG: Có viewpoint/model tương ứng và truy vết tới quyết định kiến trúc. \\ \addlinespace[1pt]
\texttt{C01.04} & QĐ 292/QĐ-BKHCN; Phụ lục, Chương 2, Mục II & A & Mô hình tham chiếu ứng dụng phải được xem xét khi mô tả thành phần ứng dụng. & AE01 & ĐÁP ỨNG: Có viewpoint/model tương ứng và truy vết tới quyết định kiến trúc. \\ \addlinespace[1pt]
\texttt{C01.05} & QĐ 292/QĐ-BKHCN; Phụ lục, Chương 2, Mục II & A & Mô hình tham chiếu công nghệ phải được xem xét khi mô tả hiện thực hóa công nghệ. & AE01; AE22 & ĐÁP ỨNG: Viewpoint V5 và VP1/VP3/VP7/VP8 phân biệt ràng buộc công nghệ với lựa chọn triển khai. \\ \addlinespace[1pt]
\texttt{C01.06} & QĐ 292/QĐ-BKHCN; Phụ lục, Chương 2, Mục II & A & Mô hình tham chiếu an toàn thông tin mạng, an ninh mạng phải được xem xét. & AE01; AE14; AE23 & ĐÁP ỨNG: Viewpoint V6 có control profile và liên kết tới interface/data classification. \\ \addlinespace[1pt]
\texttt{C02.01} & QĐ 2439/QĐ-TTg; Điều 1; Khung ban hành kèm theo, Chương I, Mục I.1 & A & Kiến trúc dữ liệu phải thống nhất nguyên tắc, cấu trúc và mô hình dữ liệu với Khung kiến trúc dữ liệu quốc gia. & AE01; AE02; AE16 & ĐÁP ỨNG: Mô hình dữ liệu sử dụng cùng khái niệm và có mapping rõ. \\ \addlinespace[1pt]
\texttt{C02.02} & QĐ 2439/QĐ-TTg; Điều 1; Khung ban hành kèm theo, Chương I, Mục I.1 & A & Phương thức kết nối/chia sẻ dữ liệu phải phù hợp Khung kiến trúc dữ liệu quốc gia. & AE11; AE12; AE20 & ĐÁP ỨNG: Mọi luồng liên cơ quan chỉ ra biên tích hợp và cơ chế điều phối. \\ \addlinespace[1pt]
\texttt{C02.03} & QĐ 2439/QĐ-TTg; Điều 1; Khung ban hành kèm theo, Chương I, Mục I.1 & A & Quản trị, quản lý dữ liệu chuyên ngành phải phù hợp Khung quản trị, quản lý dữ liệu quốc gia. & AE02; AE15; AE16 & ĐÁP ỨNG: Mỗi nhóm dữ liệu có owner/source/quality/lifecycle. \\ \addlinespace[1pt]
\texttt{C02.04} & QĐ 2439/QĐ-TTg; Điều 1; Khung ban hành kèm theo, Chương I, Mục I.1 & A & Từ điển dữ liệu chuyên ngành phải kế thừa/ánh xạ với Từ điển dữ liệu dùng chung. & AE02; AE18 & MỘT PHẦN: Cần schema/từ điển chính thức hoặc mẫu triển khai để xác nhận mapping field-level. \\ \addlinespace[1pt]
\texttt{C03.01} & Luật 93/2025/QH15; NĐ 268/2025/NĐ-CP; Luật 93/2025/QH15, khoản 3 Điều 20 & A & Nhà nước xây dựng, vận hành, duy trì Nền tảng số quản lý KH,CN\&ĐMST quốc gia tập trung, thống nhất và kết nối các chủ thể liên quan. & AE01; AE11; AE12; AE20 & ĐÁP ỨNG: Kiến trúc có biên kết nối quốc gia và không tạo silo độc lập. \\ \addlinespace[1pt]
\texttt{C03.02} & Luật 93/2025/QH15; NĐ 268/2025/NĐ-CP; NĐ 268, khoản 1 Điều 35 & A & Đối tượng DNKNST là doanh nghiệp được thành lập theo Luật Doanh nghiệp. & AE03; AE16 & ĐÁP ỨNG: Hồ sơ DNKNST bắt buộc có khóa liên kết doanh nghiệp hợp lệ. \\ \addlinespace[1pt]
\texttt{C03.03} & Luật 93/2025/QH15; NĐ 268/2025/NĐ-CP; NĐ 268, khoản 2 Điều 35 & A & DNKNST phải có dự án đã/cam kết được đầu tư hoặc được/cam kết hỗ trợ bởi các chủ thể được quy định. & AE04; AE17; AE19 & ĐÁP ỨNG: Có đối tượng quan hệ đầu tư/hỗ trợ, nguồn và tài liệu chứng minh. \\ \addlinespace[1pt]
\texttt{C03.04} & Luật 93/2025/QH15; NĐ 268/2025/NĐ-CP; NĐ 268, khoản 3 Điều 35 & A & Khả năng tăng trưởng doanh thu đạt tối thiểu 20\% trong 02 năm liên tiếp theo phân tích quy định. & AE04; AE15 & ĐÁP ỨNG: Có trường/artefact bằng chứng tăng trưởng và rule kiểm tra 20\%/02 năm. \\ \addlinespace[1pt]
\texttt{C03.05} & Luật 93/2025/QH15; NĐ 268/2025/NĐ-CP; NĐ 268, khoản 1 Điều 38 & A & Người/tổ chức yêu cầu công nhận tự kê khai và chịu trách nhiệm về tính hợp pháp, trung thực, chính xác của hồ sơ. & AE04; AE15 & ĐÁP ỨNG: Mỗi bản khai truy được tới chủ thể, thời điểm và phiên bản đã nộp. \\ \addlinespace[1pt]
\texttt{C03.06} & Luật 93/2025/QH15; NĐ 268/2025/NĐ-CP; NĐ 268, khoản 2 Điều 39 & A & Giấy công nhận bản giấy và bản điện tử có giá trị pháp lý như nhau. & AE06 & ĐÁP ỨNG: Hai dạng giấy tham chiếu cùng record/decision và trạng thái. \\ \addlinespace[1pt]
\texttt{C03.07} & Luật 93/2025/QH15; NĐ 268/2025/NĐ-CP; NĐ 268, khoản 3 Điều 39 & A & Giấy công nhận có thời hạn 05 năm từ ngày cấp. & AE06; AE10 & ĐÁP ỨNG: Ngày hết hạn = ngày cấp + 05 năm, có lịch sử thay đổi. \\ \addlinespace[1pt]
\texttt{C03.08} & Luật 93/2025/QH15; NĐ 268/2025/NĐ-CP; NĐ 268, khoản 4 Điều 40 & A & UBND cấp tỉnh có thẩm quyền công nhận DNKNST trong phạm vi quản lý. & AE04; AE20 & ĐÁP ỨNG: Hồ sơ chỉ được quyết định bởi chủ thể/cơ quan đúng phạm vi. \\ \addlinespace[1pt]
\texttt{C03.09} & Luật 93/2025/QH15; NĐ 268/2025/NĐ-CP; NĐ 268, khoản 5 Điều 40 & A & Cơ quan có thẩm quyền chịu trách nhiệm cấp, cấp lại, chấm dứt và hủy bỏ hiệu lực Giấy công nhận. & AE06; AE20 & ĐÁP ỨNG: Mọi transition pháp lý có cơ quan, quyết định và căn cứ. \\ \addlinespace[1pt]
\texttt{C03.10} & Luật 93/2025/QH15; NĐ 268/2025/NĐ-CP; NĐ 268, điểm d khoản 1 Điều 41 & A & Hội đồng tư vấn cho DNKNST có 05--07 thành viên. & AE05; AE18 & ĐÁP ỨNG: Không cho chốt hội đồng DNKNST ngoài khoảng 05--07. \\ \addlinespace[1pt]
\texttt{C03.11} & Luật 93/2025/QH15; NĐ 268/2025/NĐ-CP; NĐ 268, khoản 2 Điều 41 & A & Hội đồng có Chủ tịch, Phó Chủ tịch, ủy viên, thư ký; 02 ủy viên phản biện; ít nhất 1/2 thành viên là chuyên gia đáp ứng kinh nghiệm. & AE05; AE18 & ĐÁP ỨNG: Rule kiểm tra đủ vai trò, 02 phản biện và tỷ lệ chuyên gia. \\ \addlinespace[1pt]
\texttt{C03.12} & Luật 93/2025/QH15; NĐ 268/2025/NĐ-CP; NĐ 268, điểm a khoản 3 Điều 41 & A & Phiên họp phải có ít nhất 2/3 tổng số thành viên, có Chủ tịch/Phó Chủ tịch và đủ 02 phản biện. & AE05; AE18 & ĐÁP ỨNG: Không cho kết luận nếu quorum/roles không đạt. \\ \addlinespace[1pt]
\texttt{C03.13} & Luật 93/2025/QH15; NĐ 268/2025/NĐ-CP; NĐ 268, điểm d khoản 4 Điều 41 & A & Kết quả đồng ý khi trên 2/3 số thành viên có mặt bỏ phiếu Đồng ý. & AE05; AE18 & ĐÁP ỨNG: Kết quả tự động/kiểm tra đúng ngưỡng >2/3. \\ \addlinespace[1pt]
\texttt{C03.14} & Luật 93/2025/QH15; NĐ 268/2025/NĐ-CP; NĐ 268, điểm đ khoản 4 Điều 41 & A & Thư ký công khai kết quả biểu quyết, lập biên bản và có chữ ký xác nhận của thành viên dự họp. & AE05; AE15 & ĐÁP ỨNG: Biên bản chứa kết quả và dấu xác nhận thành viên. \\ \addlinespace[1pt]
\texttt{C03.15} & Luật 93/2025/QH15; NĐ 268/2025/NĐ-CP; NĐ 268, khoản 1 Điều 43 & A & Hồ sơ công nhận được gửi trực tiếp hoặc trực tuyến tới cơ quan có thẩm quyền. & AE04; AE07; AE24; AE26 & ĐÁP ỨNG: Hồ sơ từ hai kênh được định danh và xử lý thống nhất. \\ \addlinespace[1pt]
\texttt{C03.16} & Luật 93/2025/QH15; NĐ 268/2025/NĐ-CP; NĐ 268, điểm a khoản 2 Điều 43 & A & Trong 05 ngày làm việc phải trả lời về tính hợp lệ hồ sơ công nhận. & AE04; AE10; AE24; AE26 & ĐÁP ỨNG: Có timestamp tiếp nhận và deadline; cảnh báo/ghi bằng chứng trả lời. \\ \addlinespace[1pt]
\texttt{C03.17} & Luật 93/2025/QH15; NĐ 268/2025/NĐ-CP; NĐ 268, điểm b khoản 2 Điều 43 & A & Trong 15 ngày từ hồ sơ hợp lệ phải công nhận/cấp giấy hoặc từ chối. & AE04; AE06; AE10; AE24; AE26 & ĐÁP ỨNG: Có deadline từ mốc hồ sơ hợp lệ đến outcome. \\ \addlinespace[1pt]
\texttt{C03.18} & Luật 93/2025/QH15; NĐ 268/2025/NĐ-CP; NĐ 268, điểm c khoản 2 Điều 43 & A & Khi cần thiết, cơ quan có thẩm quyền có thể thành lập Hội đồng tư vấn công nhận. & AE04; AE05; AE24; AE26 & ĐÁP ỨNG: Có transition tạo hội đồng nhưng không bắt buộc mọi hồ sơ. \\ \addlinespace[1pt]
\texttt{C03.19} & Luật 93/2025/QH15; NĐ 268/2025/NĐ-CP; NĐ 268, điểm c khoản 2 Điều 43 & A & Nếu cần giải trình theo Hội đồng, trong 05 ngày làm việc sau họp phải yêu cầu giải trình/hoàn thiện. & AE04; AE05; AE10; AE24; AE26 & ĐÁP ỨNG: Deadline 05 ngày tính từ timestamp phiên họp. \\ \addlinespace[1pt]
\texttt{C03.20} & Luật 93/2025/QH15; NĐ 268/2025/NĐ-CP; NĐ 268, điểm c khoản 2 Điều 43 & A & Người nộp có 15 ngày để giải trình/hoàn thiện; quá hạn thì dừng xử lý. & AE04; AE10; AE24; AE26 & ĐÁP ỨNG: Quá 15 ngày không bổ sung => trạng thái dừng có căn cứ. \\ \addlinespace[1pt]
\texttt{C03.21} & Luật 93/2025/QH15; NĐ 268/2025/NĐ-CP; NĐ 268, điểm c khoản 2 Điều 43 & A & Nếu đồng ý, trong 07 ngày làm việc từ họp Hội đồng hoặc nhận hồ sơ hoàn thiện phải ra quyết định/cấp giấy. & AE04; AE06; AE10; AE24; AE26 & ĐÁP ỨNG: Deadline 07 ngày theo đúng mốc khởi tạo. \\ \addlinespace[1pt]
\texttt{C03.22} & Luật 93/2025/QH15; NĐ 268/2025/NĐ-CP; NĐ 268, điểm c khoản 2 Điều 43 & A & Nếu không đồng ý, trong 07 ngày làm việc phải thông báo kết quả và nêu rõ lý do. & AE04; AE10; AE24; AE26 & ĐÁP ỨNG: Có reason code/text, timestamp và bằng chứng gửi. \\ \addlinespace[1pt]
\texttt{C03.23} & Luật 93/2025/QH15; NĐ 268/2025/NĐ-CP; NĐ 268, khoản 3 Điều 43 & A & Nếu muốn tiếp tục công nhận, phải thực hiện thủ tục trước khi giấy hết hạn 60 ngày. & AE06; AE10; AE24 & ĐÁP ỨNG: Cảnh báo/khả năng mở hồ sơ tiếp tục trước 60 ngày. \\ \addlinespace[1pt]
\texttt{C03.24} & Luật 93/2025/QH15; NĐ 268/2025/NĐ-CP; NĐ 268, điểm a khoản 3 Điều 44 & A & Trong 05 ngày làm việc phải trả lời tính hợp lệ hồ sơ cấp lại. & AE06; AE10; AE24 & ĐÁP ỨNG: Deadline và bằng chứng trả lời được lưu. \\ \addlinespace[1pt]
\texttt{C03.25} & Luật 93/2025/QH15; NĐ 268/2025/NĐ-CP; NĐ 268, điểm b khoản 3 Điều 44 & A & Trong 15 ngày từ hồ sơ hợp lệ phải cấp lại hoặc từ chối. & AE06; AE10; AE24 & ĐÁP ỨNG: Có deadline và outcome cấp lại/từ chối. \\ \addlinespace[1pt]
\texttt{C03.26} & Luật 93/2025/QH15; NĐ 268/2025/NĐ-CP; NĐ 268, khoản 4 Điều 44 & A & Giấy cấp lại giữ thời hạn hiệu lực của giấy đã cấp. & AE06; AE24 & ĐÁP ỨNG: Expiry cấp lại không bị reset. \\ \addlinespace[1pt]
\texttt{C03.27} & Luật 93/2025/QH15; NĐ 268/2025/NĐ-CP; NĐ 268, khoản 1 Điều 45 & A & Các căn cứ chấm dứt hiệu lực Giấy công nhận phải được phản ánh trong quyết định trạng thái. & AE06; AE18; AE24; AE26 & ĐÁP ỨNG: Transition chấm dứt bắt buộc có ground hợp lệ. \\ \addlinespace[1pt]
\texttt{C03.28} & Luật 93/2025/QH15; NĐ 268/2025/NĐ-CP; NĐ 268, khoản 2 Điều 45 & A & Các căn cứ hủy bỏ hiệu lực Giấy công nhận phải được phản ánh trong quyết định trạng thái. & AE06; AE18; AE24; AE26 & ĐÁP ỨNG: Transition hủy bỏ bắt buộc có ground hợp lệ. \\ \addlinespace[1pt]
\texttt{C03.29} & Luật 93/2025/QH15; NĐ 268/2025/NĐ-CP; NĐ 268, điểm a khoản 3 Điều 45 & A & Trong 05 ngày làm việc từ yêu cầu chấm dứt/hủy bỏ phải trả lời tính hợp lệ. & AE06; AE10; AE24; AE26 & ĐÁP ỨNG: Deadline và bằng chứng trả lời được lưu. \\ \addlinespace[1pt]
\texttt{C03.30} & Luật 93/2025/QH15; NĐ 268/2025/NĐ-CP; NĐ 268, điểm b khoản 3 Điều 45 & A & Trong 15 ngày từ yêu cầu hợp lệ phải quyết định chấm dứt/hủy/từ chối; có thể lập hội đồng. & AE05; AE06; AE10; AE24; AE26 & ĐÁP ỨNG: Có deadline, quyết định và audit trail. \\ \addlinespace[1pt]
\texttt{C03.31} & Luật 93/2025/QH15; NĐ 268/2025/NĐ-CP; NĐ 268, khoản 1 Điều 48 & A & UBND cấp tỉnh nơi doanh nghiệp đặt trụ sở chính có thẩm quyền tiếp nhận, thẩm định, cấp/thay đổi/cấp lại/thu hồi/hủy Giấy chứng nhận DNKH\&CN. & AE07; AE08; AE09; AE20; AE24 & ĐÁP ỨNG: Ca kiểm thử phải chứng minh hồ sơ được định tuyến đúng tỉnh và quyết định được ghi tại AE24, không tại lớp trung ương. \\ \addlinespace[1pt]
\texttt{C03.32} & Luật 93/2025/QH15; NĐ 268/2025/NĐ-CP; NĐ 268, khoản 3 Điều 48 & A & UBND cấp tỉnh định kỳ cập nhật dữ liệu vòng đời chứng nhận DNKH\&CN lên nền tảng quốc gia. & AE06; AE10; AE11; AE12; AE15; AE24; AE25 & ĐÁP ỨNG: Mỗi thay đổi tại AE24 tạo bản ghi đồng bộ hoặc truy vấn được ở AE25 mà không thay nguồn pháp lý. \\ \addlinespace[1pt]
\texttt{C03.33} & Luật 93/2025/QH15; NĐ 268/2025/NĐ-CP; NĐ 268, khoản 1 Điều 50 & A & Hồ sơ cấp Giấy chứng nhận DNKH\&CN được nộp qua hệ thống dịch vụ công. & AE07; AE20; AE24 & ĐÁP ỨNG: Hình ranh giới thể hiện hệ thống TTHC bên ngoài và contract hồ sơ/kết quả. \\ \addlinespace[1pt]
\texttt{C03.34} & Luật 93/2025/QH15; NĐ 268/2025/NĐ-CP; NĐ 268, điểm b khoản 2 Điều 50 & A & Hồ sơ có thể dùng các văn bản công nhận/kết quả từ chuyển giao công nghệ, sở hữu trí tuệ và kết quả nhiệm vụ KH,CN\&ĐMST. & AE19; AE15; AE16; AE24 & ĐÁP ỨNG: Mỗi evidence record lưu loại, nguồn, định danh và trạng thái xác minh. \\ \addlinespace[1pt]
\texttt{C03.35} & Luật 93/2025/QH15; NĐ 268/2025/NĐ-CP; NĐ 268, khoản 1 Điều 51 & A & Trong 03 ngày làm việc từ tiếp nhận hồ sơ DNKH\&CN phải kiểm tra đầy đủ/hợp lệ. & AE08; AE10; AE24; AE26 & ĐÁP ỨNG: Có timestamp và deadline 03 ngày. \\ \addlinespace[1pt]
\texttt{C03.36} & Luật 93/2025/QH15; NĐ 268/2025/NĐ-CP; NĐ 268, khoản 2 Điều 51 & A & Trong 15 ngày từ hồ sơ hợp lệ phải cấp hoặc từ chối cấp Giấy chứng nhận. & AE08; AE06; AE10; AE24; AE26 & ĐÁP ỨNG: Có deadline và evidence của quyết định. \\ \addlinespace[1pt]
\texttt{C03.37} & Luật 93/2025/QH15; NĐ 268/2025/NĐ-CP; NĐ 268, khoản 3 Điều 51 & A & Thẩm định có thể theo tự thẩm định, chuyên gia độc lập hoặc Hội đồng. & AE08; AE05; AE24; AE26 & ĐÁP ỨNG: Mỗi hồ sơ ghi mode thẩm định và artefact đánh giá tương ứng. \\ \addlinespace[1pt]
\texttt{C03.38} & Luật 93/2025/QH15; NĐ 268/2025/NĐ-CP; NĐ 268, khoản 3 Điều 52 & A & Trong 10 ngày từ hồ sơ hợp lệ phải cấp/từ chối thay đổi nội dung hoặc cấp lại Giấy chứng nhận DNKH\&CN. & AE06; AE08; AE10; AE24; AE26 & ĐÁP ỨNG: Deadline và outcome gắn version mới. \\ \addlinespace[1pt]
\texttt{C03.39} & Luật 93/2025/QH15; NĐ 268/2025/NĐ-CP; NĐ 268, khoản 1 Điều 53 & A & Định kỳ 03 năm từ ngày cấp phải kiểm tra điều kiện duy trì hiệu lực. & AE09; AE10; AE24; AE26 & ĐÁP ỨNG: Tạo được lịch hậu kiểm từ ngày cấp và lưu kết quả. \\ \addlinespace[1pt]
\texttt{C03.40} & Luật 93/2025/QH15; NĐ 268/2025/NĐ-CP; NĐ 268, khoản 2 Điều 53 & A & Thu hồi Giấy chứng nhận chỉ theo các căn cứ được quy định. & AE09; AE18; AE24; AE26 & ĐÁP ỨNG: Transition thu hồi có ground hợp lệ và evidence. \\ \addlinespace[1pt]
\texttt{C03.41} & Luật 93/2025/QH15; NĐ 268/2025/NĐ-CP; NĐ 268, khoản 3 Điều 53 & A & Hủy bỏ hiệu lực Giấy chứng nhận chỉ theo các căn cứ được quy định. & AE09; AE18; AE24; AE26 & ĐÁP ỨNG: Transition hủy bỏ có ground hợp lệ và evidence. \\ \addlinespace[1pt]
\texttt{C03.42} & Luật 93/2025/QH15; NĐ 268/2025/NĐ-CP; NĐ 268, khoản 4 Điều 53 & A & Cơ quan có thẩm quyền có thể lập Hội đồng tư vấn khi xem xét thu hồi/hủy bỏ. & AE09; AE05; AE24; AE26 & ĐÁP ỨNG: Có thể tạo hội đồng và liên kết kết quả với quyết định. \\ \addlinespace[1pt]
\texttt{C03.43} & Luật 93/2025/QH15; NĐ 268/2025/NĐ-CP; NĐ 268, khoản 5 Điều 53 & A & Việc thu hồi/hủy bỏ phải công khai trên cổng thông tin của cơ quan ít nhất 30 ngày kể từ quyết định. & AE09; AE21; AE10; AE24; AE26 & ĐÁP ỨNG: Có mốc bắt đầu, thời hạn >=30 ngày và bằng chứng URL/ấn phẩm. \\ \addlinespace[1pt]
\texttt{C03.44} & Luật 93/2025/QH15; NĐ 268/2025/NĐ-CP; NĐ 268, khoản 2 Điều 54 & A & DNKH\&CN báo cáo đầy đủ/chính xác trên nền tảng trực tuyến của Bộ trước 15/12 hằng năm. & AE10; AE15; AE24 & ĐÁP ỨNG: Mỗi kỳ có trạng thái nộp, timestamp và cảnh báo trước 15/12. \\ \addlinespace[1pt]
\texttt{C03.45} & Luật 93/2025/QH15; NĐ 268/2025/NĐ-CP; NĐ 268, khoản 4 Điều 65 & A & Bộ KH\&CN xây dựng, quản lý, phát triển, duy trì mạng lưới/CSDL và thiết lập kết nối, liên thông, chia sẻ thông tin. & AE11; AE12; AE17; AE20; AE22; AE24 & ĐÁP ỨNG: Responsibility model chỉ rõ ai quản trị CSDL/liên thông và platform không chiếm thẩm quyền nghiệp vụ. \\ \addlinespace[1pt]
\texttt{C03.46} & Luật 93/2025/QH15; NĐ 268/2025/NĐ-CP; NĐ 268, khoản 2 Điều 66 & A & UBND cấp tỉnh/cơ quan được ủy quyền thông báo Bộ trong 15 ngày từ quyết định công nhận/cấp lại/chấm dứt/hủy bỏ và Giấy chứng nhận DNKH\&CN. & AE10; AE11; AE12; AE15; AE24; AE25 & ĐÁP ỨNG: Giả lập central unavailable: địa phương vẫn ghi quyết định; hàng đợi/outbox giữ event và gửi lại trong giới hạn nghiệp vụ. \\ \addlinespace[1pt]
\texttt{C03.47} & Luật 93/2025/QH15; NĐ 268/2025/NĐ-CP; NĐ 268, khoản 4 Điều 66 & A & Cấp tỉnh cung cấp/cập nhật thông tin mạng lưới, CSDL và công bố thông tin hoạt động ĐMST/KNST trên cổng địa phương. & AE17; AE20; AE21; AE24 & ĐÁP ỨNG: Có responsibility boundary và contract tối thiểu cho dữ liệu công bố địa phương. \\ \addlinespace[1pt]
\texttt{C03.48} & Luật 93/2025/QH15; NĐ 268/2025/NĐ-CP; NĐ 268, khoản 5 Điều 66 & A & Cấp tỉnh theo dõi/tổng hợp/đánh giá và báo cáo Bộ trước 31/3 hằng năm. & AE10; AE17; AE15; AE24; AE25 & ĐÁP ỨNG: Báo cáo trung ương truy được dữ liệu từng tỉnh, kỳ, nguồn và thời điểm đồng bộ. \\ \addlinespace[1pt]
\texttt{C04.01} & Luật Dữ liệu 60/2024/QH15; điểm a khoản 2 Điều 11 & A & Cơ quan nhà nước dùng thống nhất bảng mã danh mục dùng chung và thống nhất với dữ liệu chủ trong CSDL quốc gia. & AE02; AE16 & ĐÁP ỨNG: Có mapping dictionary/key tới nguồn quốc gia. \\ \addlinespace[1pt]
\texttt{C04.02} & Luật Dữ liệu 60/2024/QH15; điểm b khoản 2 Điều 11 & A & Dữ liệu đã có trong CSDL được kết nối, chia sẻ thì không thu thập lại. & AE03; AE11; AE16 & ĐÁP ỨNG: Trường dữ liệu nguồn ngoài được đánh dấu source và không yêu cầu nhập lại nếu truy cập được. \\ \addlinespace[1pt]
\texttt{C04.03} & Luật Dữ liệu 60/2024/QH15; điểm c khoản 2 Điều 11 & A & Dữ liệu phục vụ quản lý nhà nước/TTHC/dịch vụ công phải được tạo lập, số hóa theo pháp luật. & AE07; AE15 & ĐÁP ỨNG: Record có nguồn tạo lập, chủ thể và loại chứng từ. \\ \addlinespace[1pt]
\texttt{C04.04} & Luật Dữ liệu 60/2024/QH15; điểm đ khoản 2 Điều 11 & A & Dữ liệu thu từ số hóa phải xác thực, truy nguyên được tới bản số hóa giấy tờ/tài liệu. & AE15 & ĐÁP ỨNG: Từ record truy ngược được artefact số hóa nguồn. \\ \addlinespace[1pt]
\texttt{C04.05} & Luật Dữ liệu 60/2024/QH15; khoản 1 Điều 12 & A & Chất lượng dữ liệu gồm chính xác, hợp lệ, toàn vẹn, đầy đủ, kịp thời và thống nhất. & AE15 & ĐÁP ỨNG: Mỗi dataset quan trọng có rule/metric cho các chiều áp dụng. \\ \addlinespace[1pt]
\texttt{C04.06} & Luật Dữ liệu 60/2024/QH15; điểm a khoản 2 Điều 12 & A & Cơ quan quản lý CSDL phải áp dụng đồng bộ tiêu chuẩn/quy chuẩn và quy trình bảo đảm chất lượng. & AE02; AE15 & ĐÁP ỨNG: Metadata chỉ ra chuẩn/rule chất lượng đang áp dụng. \\ \addlinespace[1pt]
\texttt{C04.07} & Luật Dữ liệu 60/2024/QH15; điểm b khoản 2 Điều 12 & A & Phải thường xuyên kiểm tra, giám sát, khắc phục, đồng bộ và cập nhật dữ liệu. & AE11; AE15 & ĐÁP ỨNG: Có log kiểm tra/sync/correction và trạng thái lỗi. \\ \addlinespace[1pt]
\texttt{C04.08} & Luật Dữ liệu 60/2024/QH15; khoản 1 Điều 13 & A & Cơ quan nhà nước phải phân loại dữ liệu theo chia sẻ, mức quan trọng và tiêu chí quản trị/bảo vệ. & AE14; AE15; AE23 & ĐÁP ỨNG: Mỗi nhóm dữ liệu có classification profile; catalog giá trị cụ thể kế thừa nguồn có thẩm quyền. \\ \addlinespace[1pt]
\texttt{C05.01} & NĐ 278/2025/NĐ-CP; khoản 1 Điều 4 & A & Chia sẻ dữ liệu phải kịp thời, đầy đủ, đúng mục đích, an toàn và đúng thẩm quyền. & AE14; AE20 & ĐÁP ỨNG: Mỗi interface có scope, purpose, authority và audit. \\ \addlinespace[1pt]
\texttt{C05.02} & NĐ 278/2025/NĐ-CP; khoản 2 Điều 4 & A & Mọi chia sẻ dữ liệu bắt buộc phải qua Nền tảng chia sẻ, điều phối dữ liệu để giám sát, truy vết, đánh giá. & AE12; AE20; AE15 & ĐÁP ỨNG: Không có luồng bắt buộc point-to-point vượt ngoài biên quy định. \\ \addlinespace[1pt]
\texttt{C05.03} & NĐ 278/2025/NĐ-CP; khoản 3 Điều 4 & A & Chia sẻ theo danh mục CSDL/dữ liệu bắt buộc, tiêu chuẩn kỹ thuật và cơ chế có thẩm quyền. & AE02; AE11; AE20 & ĐÁP ỨNG: Mỗi interface chỉ ra dataset và chuẩn/cơ chế áp dụng. \\ \addlinespace[1pt]
\texttt{C05.04} & NĐ 278/2025/NĐ-CP; khoản 4 Điều 4 & A & Kết nối/chia sẻ phải tuân thủ Khung dữ liệu quốc gia, quản trị và từ điển dùng chung. & AE01; AE02; AE20 & ĐÁP ỨNG: Có mapping architecture/data dictionary. \\ \addlinespace[1pt]
\texttt{C05.05} & NĐ 278/2025/NĐ-CP; khoản 5 Điều 4 & A & Cơ quan quản lý CSDL phải công bố mô tả, phạm vi, chất lượng, cơ quan quản lý và điều kiện tiếp cận dữ liệu. & AE02; AE21 & ĐÁP ỨNG: Dataset catalog có đủ năm trường bắt buộc và publication mechanism. \\ \addlinespace[1pt]
\texttt{C05.06} & NĐ 278/2025/NĐ-CP; khoản 6 Điều 4 & A & Chia sẻ không được xâm phạm quyền riêng tư/bí mật cá nhân/gia đình trừ trường hợp luật quy định. & AE14 & ĐÁP ỨNG: Interface lọc dữ liệu theo quyền và purpose. \\ \addlinespace[1pt]
\texttt{C05.07} & NĐ 278/2025/NĐ-CP; khoản 7 Điều 4 & A & Dữ liệu chia sẻ phải cập nhật, chính xác. & AE11; AE15 & ĐÁP ỨNG: Payload/record có version/timestamp/source; quality checks. \\ \addlinespace[1pt]
\texttt{C05.08} & NĐ 278/2025/NĐ-CP; khoản 1 Điều 5 & A & Tích hợp, đồng bộ và sử dụng dữ liệu chủ quốc gia là bắt buộc khi xây/cập nhật/vận hành CSDL, HTTT. & AE03; AE16 & ĐÁP ỨNG: EnterpriseIdentity có upstream national master/source. \\ \addlinespace[1pt]
\texttt{C05.09} & NĐ 278/2025/NĐ-CP; khoản 2 Điều 5 & A & Một dữ liệu chủ quốc gia chỉ có một nguồn tin cậy duy nhất. & AE16; AE15 & ĐÁP ỨNG: Mỗi thuộc tính master có exactly one authoritative source role. \\ \addlinespace[1pt]
\texttt{C05.10} & NĐ 278/2025/NĐ-CP; khoản 3 Điều 5 & A & Dữ liệu chủ được thiết lập từ mã khóa định danh và danh mục dữ liệu chủ quốc gia. & AE02; AE03; AE16 & ĐÁP ỨNG: Mỗi master entity có identifier mapping documented. \\ \addlinespace[1pt]
\texttt{C05.11} & NĐ 278/2025/NĐ-CP; khoản 1 Điều 6 & A & Bộ/cơ quan trung ương xác định, công bố, cập nhật dữ liệu chủ chuyên ngành và tích hợp vào Từ điển dữ liệu dùng chung. & AE02; AE16; AE18 & ĐÁP ỨNG: Danh mục master chuyên ngành có owner/version và mapping. \\ \addlinespace[1pt]
\texttt{C05.12} & NĐ 278/2025/NĐ-CP; khoản 2 Điều 6 & A & Dữ liệu chủ chuyên ngành phải có khả năng mở rộng, tích hợp, liên thông, truy xuất với dữ liệu chủ quốc gia. & AE03; AE11; AE16 & ĐÁP ỨNG: Từ domain entity truy xuất được national identifier/source. \\ \addlinespace[1pt]
\texttt{C05.13} & NĐ 278/2025/NĐ-CP; điểm a khoản 2 Điều 7 & A & Truy vấn bắt buộc thực hiện qua Nền tảng chia sẻ, điều phối dữ liệu; nền tảng xác thực và phân quyền trao đổi. & AE12; AE14; AE20; AE24; AE25 & ĐÁP ỨNG: Sequence/interface chỉ ra platform-mediated query và authz. \\ \addlinespace[1pt]
\texttt{C05.14} & NĐ 278/2025/NĐ-CP; điểm b khoản 2 Điều 7 & A & Có thể đồng bộ một phần/toàn bộ dữ liệu qua Nền tảng chia sẻ, điều phối dữ liệu. & AE11; AE12; AE15; AE24; AE25 & ĐÁP ỨNG: Thiết kế chỉ rõ dataset nào query, partial sync hay full sync; mọi bản sao giữ provenance/version. \\ \addlinespace[1pt]
\texttt{C05.15} & NĐ 278/2025/NĐ-CP; khoản 3 Điều 7 & A & Agent Node là thành phần bảo mật điểm kết nối bắt buộc cho kết nối trực tiếp tới Nền tảng chia sẻ, điều phối dữ liệu. & AE12; AE20; AE24; AE25 & ĐÁP ỨNG: Hình ranh giới có Agent Node giữa Bộ và nền tảng quốc gia. \\ \addlinespace[1pt]
\texttt{C05.16} & NĐ 278/2025/NĐ-CP; khoản 2 Điều 10 & A & Dữ liệu bắt buộc được chia sẻ qua Nền tảng chia sẻ, điều phối; kết nối trực tiếp chỉ theo điều kiện/thỏa thuận quy định. & AE12; AE20; AE24; AE25 & ĐÁP ỨNG: Không có direct integration mặc định. \\ \addlinespace[1pt]
\texttt{C05.17} & NĐ 278/2025/NĐ-CP; khoản 2 Điều 23 & A & NĐ 47/2020/NĐ-CP hết hiệu lực từ khi NĐ 278 có hiệu lực. & AE01; AE18 & ĐÁP ỨNG: Corpus/reference audit không coi NĐ47 là normative current source. \\ \addlinespace[1pt]
\texttt{C05.18} & NĐ 278/2025/NĐ-CP; khoản 2--3 Điều 24 & A & Hệ thống phải hoàn tất chuẩn hóa và việc chia sẻ bắt buộc thống nhất qua nền tảng theo mốc 31/12/2026. & AE11; AE12; AE20; AE22 & MỘT PHẦN: Architecture supports target path; evidence of delivery milestone requires implementation plan/status. \\ \addlinespace[1pt]
\texttt{C06.01} & QĐ 1973/QĐ-BKHCN; Mục II.1; nguyên tắc kế thừa khung quốc gia & A & Khung dữ liệu Bộ kế thừa/tuân thủ kiến trúc dữ liệu quốc gia và từ điển dùng chung. & AE01; AE02 & ĐÁP ỨNG: Data view chỉ ra inheritance/mapping. \\ \addlinespace[1pt]
\texttt{C06.02} & QĐ 1973/QĐ-BKHCN; Mục II.1(d--e) & A & Các thành phần phải kết nối/liên thông/tích hợp và tránh đầu tư/lưu trữ dữ liệu trùng lặp. & AE03; AE11; AE16; AE20 & ĐÁP ỨNG: Không có competing source or duplicate point-to-point silo. \\ \addlinespace[1pt]
\texttt{C06.03} & QĐ 1973/QĐ-BKHCN; Mục II.2; nguyên tắc tạo lập dữ liệu & A & Dữ liệu được thu thập/tạo lập một lần tại cơ quan/hệ thống có thẩm quyền tạo ra dữ liệu. & AE15; AE16 & ĐÁP ỨNG: Mỗi nhóm thuộc tính có create/update authority. \\ \addlinespace[1pt]
\texttt{C06.04} & QĐ 1973/QĐ-BKHCN; Mục II.2; minh bạch chia sẻ dữ liệu & A & Dữ liệu chia sẻ phải có thông tin nguồn gốc, nội dung và phạm vi. & AE02; AE15; AE20 & ĐÁP ỨNG: Dataset catalog có source/content/scope. \\ \addlinespace[1pt]
\texttt{C06.05} & QĐ 1973/QĐ-BKHCN; Mục II.2; xác định nguồn tạo lập & A & Phải xác định điểm tạo lập/nguồn gốc cho từng nhóm dữ liệu và trách nhiệm pháp lý của nguồn. & AE16; AE15 & ĐÁP ỨNG: Authority matrix không gán một nguồn vật lý cho toàn entity. \\ \addlinespace[1pt]
\texttt{C06.06} & QĐ 1973/QĐ-BKHCN; Mục II.2; chuẩn hóa metadata/identifier & A & Dữ liệu dùng chung phải dùng chuẩn, metadata và định danh thống nhất; chia sẻ qua nền tảng tích hợp của Bộ. & AE02; AE11; AE16 & ĐÁP ỨNG: Có identifier mapping và shared integration path. \\ \addlinespace[1pt]
\texttt{C06.07} & QĐ 1973/QĐ-BKHCN; Mục III; nguyên tắc Zero Trust & A & Các kết nối phải được xác thực/ủy quyền theo nguyên tắc an toàn được khung cấp Bộ đặt ra. & AE13; AE14; AE20; AE23 & ĐÁP ỨNG: Mỗi external interface có trust boundary, authn/authz, audit/purpose concern; control implementation verified later. \\ \addlinespace[1pt]
\texttt{C06.08} & QĐ 1973/QĐ-BKHCN; Mục III; Data Lineage & A & Cần theo dõi nguồn gốc và biến đổi dữ liệu qua Data Lineage. & AE15 & ĐÁP ỨNG: Từ derived record truy được input/source/transformation. \\ \addlinespace[1pt]
\texttt{C06.09} & QĐ 1973/QĐ-BKHCN; Mục III.1.4 & A & LGSP/LDOP và API Gateway là lớp tích hợp/chia sẻ, điều phối dữ liệu cấp Bộ. & AE11; AE20; AE24; AE25 & ĐÁP ỨNG: Hình kiến trúc chỉ ra layer và responsibility. \\ \addlinespace[1pt]
\texttt{C06.10} & QĐ 1973/QĐ-BKHCN; Mục III.1.4 & A & Agent Node là điểm kết nối bảo mật giữa Bộ với các nền tảng quốc gia/Trung tâm dữ liệu quốc gia. & AE12; AE20; AE24; AE25 & ĐÁP ỨNG: System context thể hiện Agent Node/national platform. \\ \addlinespace[1pt]
\texttt{C06.11} & QĐ 1973/QĐ-BKHCN; Mục III.3.2/metadata catalog & A & Metadata phải mô tả nguồn gốc, cấu trúc, chất lượng và điều kiện khai thác dữ liệu. & AE02; AE15 & ĐÁP ỨNG: Catalog schema có origin/structure/quality/access conditions. \\ \addlinespace[1pt]
\texttt{C06.12} & QĐ 1973/QĐ-BKHCN; Mục III.3.3, STT 3 & A & Dữ liệu tổ chức/doanh nghiệp có nguồn từ CSDL quốc gia về doanh nghiệp/CSDL dùng chung Bộ Tài chính. & AE03; AE16 & ĐÁP ỨNG: Tên/định danh/status đăng ký không được platform tự tạo authoritative. \\ \addlinespace[1pt]
\texttt{C06.13} & QĐ 1973/QĐ-BKHCN; Mục III.3.3 và mô tả nền tảng dữ liệu dùng chung & A & Nền tảng dữ liệu dùng chung tích hợp dữ liệu từ hệ chuyên ngành/nguồn khác, không đồng nghĩa thay thế mọi hệ nguồn. & AE11; AE15; AE16; AE20 & ĐÁP ỨNG: Derived/shared copy luôn giữ upstream source/provenance. \\ \addlinespace[1pt]
\texttt{C06.14} & QĐ 1973/QĐ-BKHCN; Phụ lục, Mục 3, STT 2--3 & A & Danh mục dữ liệu chủ cấp Bộ bao gồm DNKH\&CN trong miền KH\&CN và DNKNST trong miền ĐMST. & AE16 & ĐÁP ỨNG: Có hai domain qualification records liên kết một enterprise identity. \\ \addlinespace[1pt]
\texttt{C07.01} & QĐ 1762/QĐ-BKHCN; Phụ lục, STT 38 & A & Danh mục CSDL chuyên ngành hiện hành liệt kê “CSDL Doanh nghiệp khởi nghiệp” do Cục Khởi nghiệp và Doanh nghiệp công nghệ chủ trì, không có một dòng riêng cho CSDL DNKH\&CN. & AE16; AE20 & CẦN LÀM RÕ: Kiến trúc phân biệt rõ kho nghiệp vụ cấp tỉnh và CSDL chuyên ngành chính thức; không tạo dòng danh mục pháp lý mới bằng quyết định kỹ thuật. \\ \addlinespace[1pt]
\texttt{C07.02} & Dự thảo Khung tiêu chuẩn dữ liệu DNKH\&CN/DNKNST (bản 2026 trong project source); Điều 4 khoản 1 & D & Dữ liệu được tạo lập một lần, cập nhật tại nguồn và sử dụng nhiều lần. & AE15; AE16 & KHÔNG CHẤM: Chỉ dùng cho chẩn đoán. \\ \addlinespace[1pt]
\texttt{C07.03} & Dự thảo Khung tiêu chuẩn dữ liệu DNKH\&CN/DNKNST (bản 2026 trong project source); Điều 4 khoản 2 & D & Mỗi dữ liệu chủ chỉ có một nguồn dữ liệu chính thức. & AE16 & KHÔNG CHẤM: Chỉ dùng cho chẩn đoán. \\ \addlinespace[1pt]
\texttt{C07.04} & Dự thảo Khung tiêu chuẩn dữ liệu DNKH\&CN/DNKNST (bản 2026 trong project source); Điều 4 khoản 3 & D & Dữ liệu phải chuẩn hóa, quản lý tập trung, thống nhất và có khả năng kết nối/chia sẻ. & AE02; AE11; AE16 & KHÔNG CHẤM: Chỉ dùng cho chẩn đoán. \\ \addlinespace[1pt]
\texttt{C07.05} & Dự thảo Khung tiêu chuẩn dữ liệu DNKH\&CN/DNKNST (bản 2026 trong project source); Điều 4 khoản 5 & D & Khuyến khích tiêu chuẩn mở và API/công nghệ hiện đại. & AE11; AE20 & KHÔNG CHẤM: Chỉ dùng cho chẩn đoán. \\ \addlinespace[1pt]
\texttt{C07.06} & Dự thảo Khung tiêu chuẩn dữ liệu DNKH\&CN/DNKNST (bản 2026 trong project source); Phụ lục I Điều 4; Phụ lục VI/VII tương ứng & D & Mã số thuế là khóa định danh chính của doanh nghiệp; mã nội bộ phải liên kết mã số thuế. & AE03; AE16 & KHÔNG CHẤM: Chỉ dùng cho chẩn đoán. \\ \addlinespace[1pt]
\texttt{C07.07} & Dự thảo Khung tiêu chuẩn dữ liệu DNKH\&CN/DNKNST (bản 2026 trong project source); Phụ lục IV, MD001 & D & MD001 DNKH\&CN gom thông tin đăng ký, hồ sơ/chứng nhận và cập nhật địa phương trong một bản ghi gốc. & AE03; AE16 & KHÔNG CHẤM: Chỉ dùng cho chẩn đoán. \\ \addlinespace[1pt]
\texttt{C07.08} & Dự thảo Khung tiêu chuẩn dữ liệu DNKH\&CN/DNKNST (bản 2026 trong project source); Phụ lục IX, MD001 & D & MD001 DNKNST gom thông tin đăng ký, hồ sơ/công nhận và cập nhật địa phương trong một bản ghi gốc. & AE03; AE16 & KHÔNG CHẤM: Chỉ dùng cho chẩn đoán. \\ \addlinespace[1pt]
\texttt{C07.09} & Dự thảo Khung tiêu chuẩn dữ liệu DNKH\&CN/DNKNST (bản 2026 trong project source); Phần căn cứ pháp lý & D & Dự thảo vẫn viện dẫn NĐ 47/2020/NĐ-CP và NĐ 13/2023/NĐ-CP. & AE01; AE18 & KHÔNG CHẤM: Chỉ dùng cho chẩn đoán. \\ \addlinespace[1pt]
\texttt{C08.01} & NĐ 69/2024/NĐ-CP; NĐ 169/2025/NĐ-CP; NĐ69, khoản 2 Điều 12 & A & Cơ quan quản lý định danh kiểm tra/xác thực thông tin tổ chức với CSDL quốc gia về đăng ký doanh nghiệp và CSDL liên quan. & AE03; AE13 & ĐÁP ỨNG: Organization identity links to official e-ID/enterprise source. \\ \addlinespace[1pt]
\texttt{C08.02} & NĐ 69/2024/NĐ-CP; NĐ 169/2025/NĐ-CP; NĐ69, khoản 3 Điều 12 & A & Kết quả đăng ký tài khoản định danh điện tử tổ chức phải được thông báo cho người thực hiện. & AE13; AE20 & ĐÁP ỨNG: System context keeps notification responsibility outside platform. \\ \addlinespace[1pt]
\texttt{C08.03} & NĐ 69/2024/NĐ-CP; NĐ 169/2025/NĐ-CP; NĐ69, điểm a khoản 4 Điều 13 & A & Cấp e-ID tổ chức không quá 03 ngày làm việc khi thông tin cần xác thực đã có trong CSDL. & AE13; AE20 & ĐÁP ỨNG: No internal API latency claim based on this deadline. \\ \addlinespace[1pt]
\texttt{C08.04} & NĐ 69/2024/NĐ-CP; NĐ 169/2025/NĐ-CP; NĐ69, điểm b khoản 4 Điều 13 & A & Cấp e-ID tổ chức không quá 15 ngày khi thông tin phải xác minh. & AE13; AE20 & ĐÁP ỨNG: Responsibility boundary documented. \\ \addlinespace[1pt]
\texttt{C08.05} & NĐ 69/2024/NĐ-CP; NĐ 169/2025/NĐ-CP; NĐ69, khoản 1 Điều 18 & A & Hệ thống của cơ quan/tổ chức kết nối e-ID phải đáp ứng ATTT tối thiểu cấp độ 3. & AE13; AE14; AE20; AE23 & MỘT PHẦN: Security architecture records level-3 precondition; operational/system-level evidence required. \\ \addlinespace[1pt]
\texttt{C08.06} & NĐ 69/2024/NĐ-CP; NĐ 169/2025/NĐ-CP; NĐ69, khoản 2 Điều 18 & A & Kết nối/chia sẻ e-ID dựa trên văn bản thống nhất nêu rõ phạm vi và mục đích. & AE13; AE20 & ĐÁP ỨNG: Contract checklist includes scope and purpose. \\ \addlinespace[1pt]
\texttt{C08.07} & NĐ 69/2024/NĐ-CP; NĐ 169/2025/NĐ-CP; NĐ69, khoản 1 Điều 19 & A & Xác thực danh tính/tài khoản điện tử qua hệ thống/nền tảng định danh và xác thực điện tử. & AE13 & ĐÁP ỨNG: Authentication flow ends at official e-ID provider. \\ \addlinespace[1pt]
\texttt{C08.08} & NĐ 69/2024/NĐ-CP; NĐ 169/2025/NĐ-CP; NĐ69, khoản 2 và khoản 4 Điều 19 & A & Cơ quan dịch vụ công có thể xác thực qua hệ e-ID/CSDL liên quan; kết quả xác thực không được tùy ý chia sẻ lại hay tái dùng làm yếu tố xác thực giao dịch khác. & AE13; AE14 & ĐÁP ỨNG: No auth-result republishing; audit purpose/consumer. \\ \addlinespace[1pt]
\texttt{C08.09} & NĐ 69/2024/NĐ-CP; NĐ 169/2025/NĐ-CP; NĐ169, khoản 1--2 Điều 43 & A & Khái niệm/phương thức xác thực mở rộng qua CSDL tổng hợp quốc gia, hệ e-ID và nền tảng dữ liệu/phương thức khác do Bộ Công an cung cấp. & AE13; AE20 & ĐÁP ỨNG: Identity interface supports replaceable authorized provider. \\ \addlinespace[1pt]
\texttt{C09.01} & Luật Bảo vệ dữ liệu cá nhân 91/2025/QH15; NĐ 356/2025/NĐ-CP; Luật 91, điểm c khoản 1 Điều 19 & A & Xử lý dữ liệu cá nhân phục vụ hoạt động cơ quan nhà nước/quản lý nhà nước có thể thuộc trường hợp không cần sự đồng ý. & AE14; AE15 & ĐÁP ỨNG: Processing record includes legal basis/purpose. \\ \addlinespace[1pt]
\texttt{C09.02} & Luật Bảo vệ dữ liệu cá nhân 91/2025/QH15; NĐ 356/2025/NĐ-CP; Luật 91, khoản 2 Điều 19 & A & Xử lý không cần consent phải có cơ chế giám sát, quy trình/trách nhiệm, biện pháp bảo vệ, đánh giá rủi ro và kiểm tra định kỳ. & AE14; AE15; AE20; AE23 & ĐÁP ỨNG: Privacy control profile có owner, purpose/legal basis, audit và risk/control-review linkage. \\ \addlinespace[1pt]
\texttt{C09.03} & Luật Bảo vệ dữ liệu cá nhân 91/2025/QH15; NĐ 356/2025/NĐ-CP; Luật 91, điểm c khoản 1 Điều 37 & A & Bên kiểm soát phải áp dụng và cập nhật biện pháp quản lý, kỹ thuật phù hợp để bảo vệ dữ liệu cá nhân. & AE14; AE18; AE23 & ĐÁP ỨNG: Security/privacy controls versioned and reviewable; implementation efficacy deferred. \\ \addlinespace[1pt]
\texttt{C09.04} & Luật Bảo vệ dữ liệu cá nhân 91/2025/QH15; NĐ 356/2025/NĐ-CP; Luật 91, điểm e khoản 1 Điều 37 & A & Bên kiểm soát phải bảo đảm các quyền của chủ thể dữ liệu cá nhân theo Điều 4. & AE14; AE15; AE23 & ĐÁP ỨNG: Architecture has explicit handling point for applicable data-subject rights; detailed workflow/legal exceptions configured per case. \\ \addlinespace[1pt]
\texttt{C09.05} & Luật Bảo vệ dữ liệu cá nhân 91/2025/QH15; NĐ 356/2025/NĐ-CP; NĐ356, khoản 2 Điều 4 & A & Khi xử lý dữ liệu cá nhân nhạy cảm, tổ chức phải thiết lập phân quyền giới hạn truy cập, quy trình xử lý và biện pháp bảo mật. & AE14; AE15 & ĐÁP ỨNG: Sensitive attributes have restricted access policy. \\ \addlinespace[1pt]
\texttt{C09.06} & Luật Bảo vệ dữ liệu cá nhân 91/2025/QH15; NĐ 356/2025/NĐ-CP; NĐ356, khoản 1--2 Điều 42 & A & NĐ356 có hiệu lực 01/01/2026 và NĐ13/2023 hết hiệu lực. & AE01; AE18 & ĐÁP ỨNG: NĐ13 not used as current normative mô hình nền. \\ \addlinespace[1pt]
\texttt{C10.01} & EIF 2017; Regulation (EU) 2024/903; EIF, Mục 3.1, p.22 & M & Quản trị liên thông phải bao quát các quyết định/thoả thuận xuyên lớp. & AE20 & THAM CHIẾU: Reference-only check; excluded from Vietnamese conformance denominator. \\ \addlinespace[1pt]
\texttt{C10.02} & EIF 2017; Regulation (EU) 2024/903; EIF, Mục 3.3, p.27 & M & Liên thông pháp lý xem xét khả năng các hệ thống hoạt động dưới các khung pháp lý tương thích. & AE20 & THAM CHIẾU: Reference-only check; excluded from Vietnamese conformance denominator. \\ \addlinespace[1pt]
\texttt{C10.03} & EIF 2017; Regulation (EU) 2024/903; EIF, Mục 3.4, p.28 & M & Liên thông tổ chức xem xét quy trình/trách nhiệm phối hợp giữa các bên. & AE20 & THAM CHIẾU: Reference-only check; excluded from Vietnamese conformance denominator. \\ \addlinespace[1pt]
\texttt{C10.04} & EIF 2017; Regulation (EU) 2024/903; EIF, Mục 3.5, p.29 & M & Liên thông ngữ nghĩa yêu cầu ý nghĩa dữ liệu được hiểu nhất quán. & AE02; AE20 & THAM CHIẾU: Reference-only check; excluded from Vietnamese conformance denominator. \\ \addlinespace[1pt]
\texttt{C10.05} & EIF 2017; Regulation (EU) 2024/903; EIF, Mục 3.6, p.30 & M & Liên thông kỹ thuật bao gồm giao diện, kết nối và hạ tầng kỹ thuật. & AE11; AE12; AE20 & THAM CHIẾU: Reference-only check; excluded from Vietnamese conformance denominator. \\ \addlinespace[1pt]
\texttt{C10.06} & EIF 2017; Regulation (EU) 2024/903; Regulation (EU) 2024/903, khoản 1--2 Điều 3; khoản 2 Điều 6 & M & Đánh giá liên thông trước thay đổi có ảnh hưởng và EIF là công cụ tham chiếu trong phạm vi EU. & AE20 & THAM CHIẾU: Reference-only check; excluded from Vietnamese conformance denominator. \\ \addlinespace[1pt]
\end{longtable}\endgroup
